\documentclass[11pt,a4paper]{article}
\usepackage[T1]{fontenc}
\usepackage{lmodern}
\usepackage[margin=27mm]{geometry}
\usepackage{amsmath,amssymb,amsthm,booktabs,microtype}
\usepackage[hidelinks]{hyperref}
\newtheorem{theorem}{Theorem}
\newtheorem{lemma}[theorem]{Lemma}
\newtheorem{proposition}[theorem]{Proposition}
\newtheorem{corollary}[theorem]{Corollary}
\theoremstyle{definition}
\newtheorem{definition}[theorem]{Definition}
\theoremstyle{remark}
\newtheorem{remark}[theorem]{Remark}
\newcommand{\Ostar}{O^{*}}
\newcommand{\E}{\mathbb E}
\newcommand{\Prb}{\mathbb P}
\emergencystretch=2em
\clubpenalty=10000
\widowpenalty=10000
\hypersetup{pdftitle={Count Guessing Makes Constraints Useful}}
\title{Count Guessing Makes Constraints Useful\\
\large Max 2-SAT for aligned blocks plus a pseudo-random remainder}
\author{}
\date{11 October 2026}
\begin{document}
\maketitle

\begin{abstract}
The note on useful constraints and group counts~\cite{RG} names two exits for
the branching algorithm: learned constraints that give useful additional
pruning, and compression of redundant structure. We combine them. If the
quadratic part of the score is a few \emph{count terms} plus a pseudo-random
remainder, the algorithm first guesses the values of the counts. Every count
term then becomes a field, the learned constraints come from the remainder
only, and early enumeration~\cite{EE,CC} prunes as for a pseudo-random
formula. The guess costs a polynomial factor. We prove a deterministic
theorem: such formulas are solved exactly in time $\Ostar(c^n)$, $c<2$, and
polynomial space, at every density, with arbitrary fields. The class contains
the pseudo-random formulas of~\cite{CC}. With probability at least $1-1/n$ it
also contains random formulas whose pairs carry independent clause parities:
with any parity bias, at every density $d\ge C\log n$, with any number of
clauses per pair and with arbitrary unit clauses
(Proposition~\ref{prop:random}). Random planted bisections and random formulas
with a bounded number of aligned groups are members as well, together with
their planted decompositions. The aligned family
of~\cite[Section~4]{RG} is its degenerate case. The simplest members need no
unknowns: the parity matrix minus its mean is pseudo-random. On strongly biased random formulas, plain early enumeration
keeps a constant fraction of $10^{-4}$ to $10^{-3}$ of all assignments
up to $n=480$; with count guessing no sampled assignment survives.
\end{abstract}

\section{The idea in clause terms}\label{sec:idea}

A learned constraint is redundant when its bad values agree with those of the
constraints already present~\cite{RG}. Agreement on a large scale has a simple
source: many clauses whose parities follow one pattern. In the aligned family
of~\cite[Section~4]{RG}, every pair carries $(u\lor v)$ or
$(\neg u\lor\neg v)$, and the score is a linear function minus
$\binom k2$, where $k$ is the number of true variables. All learned constraints
then push the same way, and the algorithm of~\cite{RG} instead sorts for every
value of $k$.

Now add noise. Let a fraction of the pairs carry the opposite parity at random.
The count term is still there, so the constraints still largely agree, and the
summed constraint stays weak. But sorting no longer solves the problem either.
The two exits combine. For a \emph{guessed} final count $k$, the count term
contributes to every single flip only a shift by a common field, which depends
on $k$ alone. Replace it by that shift. The learned constraints then come from
the noise, which is pseudo-random, and they prune as in~\cite{CC}. One of the
$n+1$ guesses is right, and in that run the optimum survives.

\section{The class}\label{sec:class}

Write the score of a 2-CNF formula on $x\in\{0,1\}^n$ as
\[
 F(x)=c_0+\sum_uh_ux_u+\sum_{u<v}J_{uv}x_ux_v .
\]
Fix the order $1,\dots,n$ and a depth $T=\lfloor\tau n\rfloor$, $0<\tau<1$.

\begin{definition}[Aligned decomposition]
Let $s_1,\dots,s_r\in\{-1,0,1\}^n$ and $\beta_1,\dots,\beta_r\in\mathbb Q$.
Put $A_{uv}=\sum_k\beta_ks_k(u)s_k(v)$ for $u\ne v$, the \emph{aligned part},
and $N=J-A$ off the diagonal, the \emph{remainder}. The \emph{counts} of $x$
are $c_k(x)=\langle s_k,x\rangle$, and $\sigma=\sum_k|\beta_k|$.
\end{definition}

If the $s_k$ are indicators of blocks, the aligned part is constant inside each
block; sums and differences of two indicators add a constant bias between two
blocks. The counts are then the group counts of~\cite{RG}, and the number of
runs below is governed by the number of distinct group counts.

\begin{definition}[Pseudo-random remainder]\label{def:prr}
Let $D=\frac1n\sum_{u\ne v}N_{uv}^2$, let $R$ be the $T\times n$ matrix with
$R_{tw}=N_{tw}$ for $w>t$ and $0$ otherwise, and $W_t=\sum_{j<t}N_{tj}^2$.
For constants $K,g>0$, the remainder is \emph{$(K,g)$-pseudo-random} if
\begin{align}
 \|R\|^2&\le K\,D,\tag{R1}\label{R1}\\
 \sum_{t\le T}\sqrt{W_t}&\ge4\sqrt3\,g\,n\sqrt D,\tag{R2}\label{R2}\\
 \sigma&\le g\sqrt D/(8\tau).\tag{R3}\label{R3}
\end{align}
\end{definition}

Condition~\eqref{R1} is the aggregate pair-count condition of~\cite{CC} for the
remainder. Condition~\eqref{R2} asks the earlier neighbours to carry weight, as
in~\cite{CA}. Condition~\eqref{R3} compares the strength of the aligned
part per pair with the scale $\sqrt D$ of the remainder; it holds as soon as
$D$ is large. No condition is placed on the fields $h$ or on the aligned part
otherwise.

\paragraph{Examples.}
\begin{enumerate}
 \item $r=0$: the pseudo-random formulas of~\cite{CC,CA}.
 \item \emph{Mean-centred}: $r=1$, $s=\mathbf 1$, and $\beta$ the mean of
 $J_{uv}$ over all pairs $u\ne v$. Nothing is unknown, and all conditions can
 be checked in polynomial time (Corollary~\ref{cor:mean}). Random formulas
 whose pairs carry independent clause parities with a common mean are members
 with probability at least $1-1/n$: with any parity bias, at every density
 $d\ge C\log n$, with any number of clauses per pair and with arbitrary unit
 clauses (Section~\ref{sec:random}, Corollary~\ref{cor:random}). Both standard
 models are of this kind: one clause per pair, and each of the four clauses of
 a pair present independently.
 \item \emph{Planted bisection}: $r=1$, $s\in\{\pm1\}^n$: pairs on the same
 side have one parity bias, pairs across the other. Random planted bisections,
 whose pairs carry independent parities with mean $\beta s(u)s(v)$, are members
 by Proposition~\ref{prop:random}(ii) when $s$ is given.
 \item \emph{Few groups}: $s_k$ the indicators of $r=O(1)$ groups and $\beta_k$
 a bias per pair of groups, written as a sum of terms of this form. Random
 formulas whose pairs carry independent parities with a mean that depends only
 on the groups of the two variables are members by
 Proposition~\ref{prop:random}(ii) when the groups are given.
 \item \emph{Degenerate}: $N=0$. Then every run below optimizes a linear
 function, which is the sorting argument of~\cite[Section~4]{RG}. The
 theorem does not need this case.
\end{enumerate}

\section{The algorithm}\label{sec:algorithm}

For every vector $c=(c_1,\dots,c_r)$ of integers with
$|c_k|\le\|s_k\|_1$ run the following, and return the best assignment found
over all runs.
\begin{enumerate}
 \item Shift the fields: $h^c=h+\sum_k\beta_kc_ks_k$, and let
 $F'_c(x)=c_0+\langle h^c,x\rangle+\sum_{u<v}N_{uv}x_ux_v$.
 \item Branch on $x_1,\dots,x_T$, visiting all $2^T$ prefixes. At step $t$
 let
 \[
  \Delta_t=h^c_t+\sum_{j<t}N_{tj}x_j+\tfrac12\sum_{w>t}N_{tw},
 \]
 the gain of $x_t=1$ over $x_t=0$ under $F'_c$ when the later variables are
 set to $\frac12$. The preferred value is $p_t=[\Delta_t\ge0]$. If $x_t\ne
 p_t$, then $t$ is a \emph{creator}, and $\gamma_t(x)$ denotes the gain of
 flipping $x_t$ to $p_t$ under $F'_c$, a linear function of the other
 variables.
 \item At each prefix, enumerate the completions satisfying the summed
 constraint
 \begin{equation}
  \sum_{t\text{ creator}}\gamma_t(x)\le\sigma\cdot\#\{\text{creators}\},
  \label{eq:summed}
 \end{equation}
 with the output-sensitive generator of~\cite{EE}, and evaluate $F$ on each.
\end{enumerate}
There are at most $(2n+1)^r$ runs. Every run uses polynomial space, and its
time is $\Ostar(2^T+W_c)$, where $W_c$ is the number of assignments $x$ that
satisfy~\eqref{eq:summed} for their own prefix~\cite[Lemma~2]{EE}.

\section{Exactness}\label{sec:exact}

\begin{lemma}[Gains]\label{lem:gain}
Let $c(x)=c$. Flipping $x_t$ from $0$ to $1$ changes $F$ and $F'_c$ equally.
Flipping $x_t$ from $1$ to $0$ changes $F$ by
$\sum_k\beta_ks_k(t)^2\ge-\sigma$ more than it changes $F'_c$.
\end{lemma}

\begin{proof}
The aligned part of $F$ is $\sum_k\frac{\beta_k}2\bigl(c_k(x)^2-\sum_u
s_k(u)^2x_u\bigr)$, since $x_u^2=x_u$. Flipping $x_t$ up changes it by
$\sum_k\beta_kc_ks_k(t)$, the same as the shift in $F'_c$. Flipping $x_t$
down changes it by $\sum_k\beta_k(s_k(t)^2-c_ks_k(t))$, while the shift
changes by $-\sum_k\beta_kc_ks_k(t)$. The remainder and $h$ contribute
equally to both.
\end{proof}

\begin{lemma}[Exactness]\label{lem:exact}
The algorithm returns an optimal assignment.
\end{lemma}

\begin{proof}
Add to $F$ a generic perturbation $\langle\theta,x\rangle$ with $\|\theta\|_1$ below
$1/q$, where $q$ is a common denominator of all coefficients. Let $x^\circ$
be the unique maximizer of $F+\langle\theta,x\rangle$; it maximizes $F$. Consider
the run $c=c(x^\circ)$, add the same perturbation to $F'_c$, and let $t$ be a
creator on the path of $x^\circ$. If the perturbed gain of flipping $x_t$ to
$p_t$ at $x^\circ$ were at least $\sigma$, Lemma~\ref{lem:gain} would make the
perturbed gain under $F$ nonnegative, contradicting uniqueness. Summing the
strict inequalities over the creators and removing the perturbation gives
$\sum_t\gamma_t(x^\circ)<\sigma\cdot\#\{\text{creators}\}+1/q$. The left side
minus $\sigma\cdot\#\{\text{creators}\}$ is a multiple of $1/q$, so it is at
most zero. Hence $x^\circ$ satisfies~\eqref{eq:summed} at its own prefix and
is enumerated.
\end{proof}

The argument needs no property of $N$: the algorithm is exact for every
decomposition, and the conditions of Definition~\ref{def:prr} only bound its
running time.

\section{Running time}\label{sec:time}

Fix a run $c$ and let $x$ be uniform; put $u=2x-1\in\{\pm1\}^n$. Write $H_t=
h^c_t+\frac12\sum_{w\ne t}N_{tw}$ and let $L$ be the $T\times T$ matrix with
$L_{tj}=N_{tj}$ for $j<t$. Then $\Delta_t=H_t+\frac12(Lu)_t$, and
$L^{\mathsf T}$ is the leading block of $R$, so $\|L\|\le\|R\|$.

\begin{lemma}[Passing]\label{lem:pass}
With $I$ the set of creators, $x$ satisfies~\eqref{eq:summed} if and only if
$S'+Q\le0$, where
\[
 S'=\sum_{t\in I}\bigl(|\Delta_t|-\sigma\bigr),\qquad
 Q=-\tfrac12\sum_{w=1}^nu_w\sum_{t\in I,\,t<w}u_tN_{tw}.
\]
\end{lemma}

\begin{proof}
For a creator, $\gamma_t(x)=(2p_t-1)\bigl(\Delta_t+\frac12\sum_{w>t}
N_{tw}u_w\bigr)$, $(2p_t-1)\Delta_t=|\Delta_t|$ and $u_t=-(2p_t-1)$.
\end{proof}

\begin{lemma}[Gap profile with arbitrary fields]\label{lem:profile}
Assume \eqref{R1}, \eqref{R2} and $n\ge256\pi\tau K/g^2$. Let
$\alpha=\max\{2,\frac94\tau\sqrt K/g\}$ and $M=\alpha\sqrt{KD}$. Then
\[
 \Prb\Bigl[\sum_{t\le T}\min\{|\Delta_t|,M\}<\tfrac34g\,n\sqrt D\Bigr]
 \le4e^{-g^2n/(16\tau K)} .
\]
\end{lemma}

\begin{proof}
Split $t\le T$ into $\mathrm{Small}=\{|H_t|\le M/2\}$ and
$\mathrm{Large}=\{|H_t|>M/2\}$. Since $\|Lu\|^2\le\|L\|^2T\le KDT$, two
deterministic facts hold.
\begin{itemize}
 \item For $t\in\mathrm{Small}$, $|\Delta_t|-M\le\frac12|(Lu)_t|-\frac M2$,
 and $(a-b)_+\le a^2/(4b)$ gives
 $\min\{|\Delta_t|,M\}\ge|\Delta_t|-(Lu)_t^2/(8M)$. Summed over Small, the
 loss is at most $KDT/(8M)=T\sqrt{KD}/(8\alpha)$.
 \item For $t\in\mathrm{Large}$ with $|(Lu)_t|\le M/2$,
 $|\Delta_t|\ge M/4$. At most $4KDT/M^2=4T/\alpha^2$ indices have
 $|(Lu)_t|>M/2$.
\end{itemize}
So $\sum_t\min\{|\Delta_t|,M\}\ge G(u)-T\sqrt{KD}/(8\alpha)+\frac
M4(|\mathrm{Large}|-4T/\alpha^2)$ with $G(u)=\sum_{t\in\mathrm{Small}}|\Delta_t|$.
Each $(Lu)_t$ is a symmetric sum of independent terms with variance $W_t$, so
$\E|\Delta_t|\ge\E|\frac12(Lu)_t|\ge\frac12\sqrt{W_t/3}$ by symmetry and the
fourth-moment bound $\E|Z|\ge\sqrt{\operatorname{Var}Z/3}$~\cite{EE}. For
$t\in\mathrm{Large}$, $\frac M4\ge\frac12\sqrt{W_t/3}$ because
$W_t\le\|L\|^2\le KD$ and $\alpha\ge2$. With~\eqref{R2},
\[
 \E G+\tfrac M4|\mathrm{Large}|\ \ge\ \tfrac12\sum_{t\le T}\sqrt{W_t/3}
 \ \ge\ 2g\,n\sqrt D ,
\]
and the two losses $T\sqrt{KD}/\alpha+T\sqrt{KD}/(8\alpha)$ are at most
$\frac12g\,n\sqrt D$ by the choice of $\alpha$. The function $G$ is convex in
$u\in[-1,1]^T$ and Lipschitz with constant
$\ell=\frac12\sqrt T\|L\|\le\frac12\sqrt{TKD}$. Talagrand's
inequality~\cite{Tal95} gives $\Prb[G\le\E G-s-8\sqrt\pi\,\ell]\le
4e^{-s^2/(16\ell^2)}$; the assumption on $n$ makes $8\sqrt\pi\,\ell\le\frac14
g\,n\sqrt D$. Take $s=\frac12g\,n\sqrt D$.
\end{proof}

\begin{proposition}[One run]\label{prop:run}
Under \eqref{R1}--\eqref{R3} and $n\ge256\pi\tau K/g^2$, a uniform $x$
satisfies~\eqref{eq:summed} for its own prefix with probability at most
$6e^{-c_*n}$, where $c_*=g^2/(64\tau\alpha^2K)$. The bound does not depend on the
fields, hence not on the run.
\end{proposition}

\begin{proof}
\emph{Shortfall.} Let $Y_t=[t\in I]\min\{|\Delta_t|,M\}\in[0,M]$. Given $u_1,\dots,u_{t-1}$,
the value $\Delta_t$ is fixed and $t$ is a creator with probability
$\frac12$. So $\E[Y_t\mid\cdot]=\frac12\min\{|\Delta_t|,M\}$ and
$\E[Y_t^2\mid\cdot]\le M^2/2$. The lower tail for nonnegative
increments~\cite{EE,Mau03} gives $\sum Y_t\ge\frac12\sum_t\min\{|\Delta_t|,M\}
-a$ except with probability $e^{-a^2/(TM^2)}$. Take $a=\frac18g\,n\sqrt D$.
By Lemma~\ref{lem:profile} and~\eqref{R3},
$S'\ge\sum Y_t-\sigma T\ge(\frac38-\frac18-\frac18)g\,n\sqrt D$ outside the
two exceptional events.

\emph{Martingale.} In the order $w=1,\dots,n$, the coefficient
$\kappa_w=\sum_{t\in I,t<w}u_tN_{tw}$ of $u_w$ in $Q$ is fixed before $u_w$
is drawn, and $\sum_w\kappa_w^2=\|R^{\mathsf T}v\|^2\le KDT$ with
$v_t=u_t[t\in I]$. The conditional $\cosh$ bound gives
$\Prb[Q\le-b]\le e^{-2b^2/(KDT)}$; take $b=\frac18g\,n\sqrt D$.

Outside the three events $S'+Q>0$. The exponents are $g^2n/(64\tau\alpha^2K)$,
$g^2n/(16\tau K)$ and $g^2n/(32\tau K)$, and the first is the smallest.
\end{proof}

\begin{theorem}[Aligned blocks plus a pseudo-random remainder]\label{thm:main}
Let $F$ be a 2-CNF formula on $n$ variables with an aligned decomposition of
rank $r$ whose remainder is $(K,g)$-pseudo-random, and let
$n\ge256\pi\tau K/g^2$. Given the decomposition, Max 2-SAT on $F$ is solved
exactly by a deterministic algorithm in polynomial space and time
\[
 \Ostar\Bigl((2n+1)^r\bigl(2^{\tau n}+2^ne^{-c_*n}\bigr)\Bigr),
 \qquad c_*=\frac{g^2}{64\tau\alpha^2K},\quad
 \alpha=\max\Bigl\{2,\frac{9\tau\sqrt K}{4g}\Bigr\}.
\]
For $r=O(1)$ this is $\Ostar(c^n)$ with $c=2^{\max\{\tau,\,1-c_*/\ln2\}}<2$,
independent of the density and of the fields.
\end{theorem}

\begin{proof}
Lemma~\ref{lem:exact} gives exactness. Each run costs
$\Ostar(2^T+W_c)$, and $W_c\le2^n\cdot6e^{-c_*n}$ by
Proposition~\ref{prop:run}.
\end{proof}

\begin{corollary}[Mean-centred formulas]\label{cor:mean}
Let $\beta$ be the mean of $J_{uv}$ over all pairs $u\ne v$ and
$N=J-\beta(\mathbf 1\mathbf 1^{\mathsf T}-I)$. If $N$ satisfies
\eqref{R1}--\eqref{R3} with $\sigma=|\beta|$, then Max 2-SAT on $F$ is solved
in polynomial space and time $\Ostar(n(2^{\tau n}+2^ne^{-c_*n}))$. All
conditions can be verified in polynomial time; there are $n+1$ runs, one for
each number of true variables.
\end{corollary}

\begin{remark}[Constants and other decompositions]
The constants are very small: for the complete parity-biased formulas with
$\tau=\frac12$, \eqref{R2} holds with $g\approx0.034$
(Table~\ref{tab:membership}). The measured behaviour is far better
(Table~\ref{tab:pass}). For a planted bisection $s$ is unknown.
When $|\beta|n\gg\|N\|$, rounding the top eigenvector of $J$ recovers $s$
except for $O(\|N\|^2/(\beta^2n))$ coordinates, by the Davis--Kahan theorem.
The misplaced coordinates change $N$ by a matrix of rank at most two and norm
$O(\|N\|)$. This recovery is sketched, not proved here.
\end{remark}

\section{Random formulas}\label{sec:random}

The class of Definition~\ref{def:prr} is deterministic. This section shows that
random formulas belong to it with probability at least $1-1/n$, with any parity
bias, at every density $d\ge C\log n$, and with arbitrary unit clauses. Recall
that $J_{uv}$ is the number of mixed clauses on the pair $\{u,v\}$ minus the
number of equal clauses on it: $(u\lor v)$ and $(\neg u\lor\neg v)$ contribute
$-1$, while $(u\lor\neg v)$ and $(\neg u\lor v)$ contribute $+1$.

\begin{definition}[Independent pairs]\label{def:indep}
A random 2-CNF formula has \emph{independent pairs} if the numbers $J_{uv}$,
$u<v$, are independent random variables. Nothing else is assumed: the unit
clauses, and which literals are negated in a clause of a given parity, may be
arbitrary and may even depend on the $J_{uv}$.
\end{definition}

\paragraph{Two models.} In both, $\varepsilon\in[-1,1]$ is the parity bias,
and $d$ is the expected number of binary clauses that contain a given variable.
\begin{enumerate}
 \item[(a)] \emph{One clause per pair.} Each pair independently receives no
 clause with probability $1-p$. Otherwise it receives exactly one clause: its
 literal on $u$ is positive or negative with probability $\frac12$ each, and
 its literal on $v$ has the same polarity with probability
 $\frac{1+\varepsilon}2$ and the opposite polarity otherwise. Here $d=p(n-1)$,
 $J_{uv}\in\{-1,0,1\}$, $\E J_{uv}=-\varepsilon p$ and
 $\operatorname{Var}J_{uv}=p(1-\varepsilon^2p)$. For $\varepsilon=0$ these are
 the usual random formulas on $G(n,p)$; for $p=1$ and $\varepsilon=1$ they
 form the aligned family of~\cite[Section~4]{RG}.
 \item[(b)] \emph{Independent clauses.} Each of the four clauses on each pair
 is present independently: $(u\lor v)$ and $(\neg u\lor\neg v)$ with
 probability $q(1+\varepsilon)$ each, $(u\lor\neg v)$ and $(\neg u\lor v)$
 with probability $q(1-\varepsilon)$ each. A pair carries between $0$ and $4$
 clauses, $d=4q(n-1)$, $|J_{uv}|\le2$ and $\E J_{uv}=-4q\varepsilon$; if
 $q(1+|\varepsilon|)\le\frac12$, then $\operatorname{Var}J_{uv}\ge2q$. For
 $\varepsilon=0$ and $q=\frac14$ every clause is present with probability
 $\frac14$.
\end{enumerate}

\begin{proposition}[Random formulas are pseudo-random]\label{prop:random}
Let $0<\tau<1$, $\rho\ge1$, $C=10^9\rho/\tau$ and $n\ge n_0=10^9\rho^2/\tau^2$.
Let $F$ have independent pairs with $|J_{uv}-\E J_{uv}|\le B$ and
$\nu\le\operatorname{Var}J_{uv}\le\rho\nu$ for all $u<v$, and take one of the
following aligned parts.
\begin{enumerate}
 \item[(i)] \emph{Mean-centred}: $\E J_{uv}=\mu$ for all pairs, and
 $A=\beta(\mathbf 1\mathbf 1^{\mathsf T}-I)$ with $\beta$ the mean of the
 $J_{uv}$, as in Corollary~\ref{cor:mean}; put $\sigma_0=|\mu|$.
 \item[(ii)] \emph{Expected}: a fixed aligned decomposition with
 $A_{uv}=\E J_{uv}$ for all $u\ne v$; put $\sigma_0=\sigma$.
\end{enumerate}
If $n\nu\ge C(B^2\log n+\rho\sigma_0^2)$, then with probability at least
$1-1/n$ the remainder $N=J-A$ is $(K,g)$-pseudo-random with $K=9\rho$ and
$g=\tau^{3/2}/(12\sqrt\rho)$. Moreover $n\ge256\pi\tau K/g^2$.
\end{proposition}

\begin{proof}
Let $X_{uv}=J_{uv}-\E J_{uv}$. These are independent and centred, with
$|X_{uv}|\le B$ and $\nu\le\E X_{uv}^2\le\rho\nu$. Let $\delta=\beta-\mu$ in
case~(i) and $\delta=0$ in case~(ii); in both cases $N_{uv}=X_{uv}-\delta$ for
$u\ne v$. Put $\eta=1/100$. The assumption gives $B\sqrt{\log n}\le\sqrt{n\nu/C}$
and $\sigma_0\le\sqrt{n\nu/(C\rho)}$. Three standard tools are used.
\begin{itemize}
 \item \emph{Bernstein's inequality}~\cite{BLM13}: if $Y_1,\dots,Y_m$ are
 independent, $Y_i-\E Y_i\le\bar c$ and $\sum_i\operatorname{Var}Y_i\le\lambda$,
 then $\Prb[\sum_i(Y_i-\E Y_i)\ge s]\le\exp(-s^2/(2\lambda+2\bar cs/3))$.
 \item \emph{Talagrand's inequality}~\cite{Tal95}, as in
 Lemma~\ref{lem:profile}: if $f$ is convex and $1$-Lipschitz in independent
 variables with values in $[-B,B]$, then $f\le\E f+8\sqrt\pi\,B+s$ except with
 probability $4e^{-s^2/(16B^2)}$.
 \item \emph{The norm bound of Bandeira and van
 Handel}~\cite[Theorem~3.1 and Corollary~3.2]{BvH16}: for a fixed $T\times n$
 matrix $(a_{tw})$ and independent random signs $\xi_{tw}$,
 \[
  \E\|(\xi_{tw}a_{tw})\|\le(1+\vartheta)\Bigl(\sigma_1+\sigma_2
  +\frac{5\max_{t,w}|a_{tw}|\sqrt{\log n}}{\sqrt{\log(1+\vartheta)}}\Bigr),
  \qquad0<\vartheta\le\tfrac12,
 \]
 where $\sigma_1$ and $\sigma_2$ are the largest Euclidean norms of a row and
 of a column of $(a_{tw})$.
\end{itemize}
Each event below fails with probability at most $n^{-3}$, and there are at most
$n+4$ of them, so they all hold with probability at least $1-1/n$.

\emph{The mean and $D$.} In case~(i), $\binom n2\delta=\sum_{u<v}X_{uv}$, and
Bernstein's inequality for both tails (with $\bar c=B$ and
$\lambda=\binom n2\rho\nu$) gives $|\delta|\le\delta_0=6\sqrt{\rho\nu\log n}/n$:
the assumption makes $\bar cs/3\le\lambda$, and the exponent is then at least
$4.5(1-\frac1n)\log n$. As $n\ge n_0$, $\delta_0\le10^{-3}\sqrt\nu$ and
$n\delta_0\le10^{-3}\sqrt{n\rho\nu}$. Bernstein's inequality for the sum of the
$X_{uv}^2$ (with $\bar c=B^2$ and $\lambda=\binom n2B^2\rho\nu$) gives
$(1-\eta)n(n-1)\nu\le\sum_{u\ne v}X_{uv}^2\le(1+\eta)n(n-1)\rho\nu$. Since
$\sum_{u\ne v}X_{uv}=n(n-1)\delta$, we have
$D=\frac1n\sum_{u\ne v}X_{uv}^2-(n-1)\delta^2$, and hence
\begin{equation}\label{eq:Dbounds}
 (1-2\eta)(n-1)\nu\le D\le(1+\eta)(n-1)\rho\nu .
\end{equation}

\emph{Condition~\eqref{R1}.} Let $R_X$ be the $T\times n$ matrix with entries
$X_{tw}$ for $w>t$ and $0$ otherwise, and $U$ the $0/1$ matrix with the same
pattern. Then $R=R_X-\delta U$ and $\|U\|\le\|U\|_F\le n$. Let $X'$ be an
independent copy of $X$ and $\hat X=X-X'$. By Jensen's inequality,
$\E\|R_X\|\le\E\|R_X-R_{X'}\|=\E\|R_{\hat X}\|$. The entries of $\hat X$ are
independent and symmetric, so $R_{\hat X}$ has the law of
$(\xi_{tw}\hat X_{tw})$. The norm bound, applied conditionally on $\hat X$ with
$\vartheta=\eta$ and $|\hat X_{tw}|\le2B$, gives
\[
 \E\|R_X\|\le(1+\eta)\bigl(\E\sigma_1(\hat X)+\E\sigma_2(\hat X)\bigr)
 +102\,B\sqrt{\log n},
\]
where $\sigma_1(\hat X)^2$ and $\sigma_2(\hat X)^2$ are the largest row sum and
the largest column sum of the $\hat X_{tw}^2$ over the pattern. A row has fewer
than $n$ entries and a column at most $T\le\tau n$, each with
$\hat X_{tw}^2\le4B^2$ and $\E\hat X_{tw}^2\le2\rho\nu$. Bernstein's inequality
(with $\bar c=4B^2$, $\lambda=8nB^2\rho\nu$ and $s=2\eta n\rho\nu$) and a union
bound over the at most $2n$ rows and columns show that all row sums are at most
$2(1+\eta)n\rho\nu$ and all column sums at most $2(\tau+\eta)n\rho\nu$, except
with probability $2n\,e^{-\eta^2n\rho\nu/(5B^2)}\le n^{-3}$. Every such sum is
at most $4nB^2$, so $\E\sigma_1(\hat X)^2\le2(1+2\eta)n\rho\nu$ and
$\E\sigma_2(\hat X)^2\le2(\tau+2\eta)n\rho\nu$. With $102/\sqrt C\le0.004$,
\[
 \E\|R_X\|\le\Bigl(1.01\sqrt2\bigl(\sqrt{1.02}+\sqrt{\tau+0.02}\bigr)
 +0.004\Bigr)\sqrt{n\rho\nu}\le2.89\sqrt{n\rho\nu}.
\]
The norm $\|R_X\|$ is a convex function of the entries $X_{tw}$ and is
$1$-Lipschitz for the Euclidean norm. Talagrand's inequality with
$s=\eta\sqrt{n\rho\nu}$ gives
$\|R_X\|\le\E\|R_X\|+8\sqrt\pi\,B+\eta\sqrt{n\rho\nu}$ except with probability
$4n^{-\eta^2C/16}\le n^{-3}$. As $8\sqrt\pi\,B\le10^{-3}\sqrt{n\rho\nu}$,
\[
 \|R\|\le\|R_X\|+n|\delta|\le2.902\sqrt{n\rho\nu},\qquad
 \|R\|^2\le8.43\,n\rho\nu\le9\rho(1-2\eta)(n-1)\nu\le KD .
\]

\emph{Condition~\eqref{R2}.} Let $t_0=\lceil\tau n/10\rceil$. For $t>t_0$ the
sum $V_t=\sum_{j<t}X_{tj}^2$ has mean at least $(t-1)\nu$, and Bernstein's
inequality for its lower tail (with $\bar c=B^2$, $\lambda=(t-1)B^2\rho\nu$ and
$s=\eta(t-1)\nu$) gives $V_t\ge(1-\eta)(t-1)\nu$ except with probability
$e^{-\eta^2\tau n\nu/(30B^2\rho)}\le n^{-3}$. Then
$\sqrt{W_t}\ge\sqrt{V_t}-|\delta|\sqrt{t-1}\ge0.993\sqrt{(t-1)\nu}$. Using
$\sum_{s=t_0}^{T-1}\sqrt s\ge\frac23\bigl((T-1)^{3/2}-t_0^{3/2}\bigr)$,
$T-1\ge\tau n-2$ and $\tau n\ge10^3$,
\[
 \sum_{t\le T}\sqrt{W_t}\ge0.993\sqrt\nu\cdot\tfrac23\bigl((T-1)^{3/2}
 -t_0^{3/2}\bigr)\ge0.635\,\tau^{3/2}n^{3/2}\sqrt\nu .
\]
By~\eqref{eq:Dbounds},
$4\sqrt3\,g\,n\sqrt D\le\frac{\sqrt3}3\sqrt{1.01}\,\tau^{3/2}n^{3/2}\sqrt\nu
\le0.581\,\tau^{3/2}n^{3/2}\sqrt\nu$.

\emph{Condition~\eqref{R3}.} By~\eqref{eq:Dbounds},
$g\sqrt D/(8\tau)\ge\sqrt\tau\sqrt{0.98(n-1)\nu}/(96\sqrt\rho)
\ge0.0102\sqrt{\tau n\nu/\rho}$. In case~(ii),
$\sigma=\sigma_0\le\sqrt{n\nu/(C\rho)}\le10^{-4}\sqrt{\tau n\nu/\rho}$. In
case~(i), $\sigma=|\beta|\le\sigma_0+\delta_0\le0.002\sqrt{\tau n\nu/\rho}$.

Finally, $256\pi\tau K/g^2=256\pi\cdot1296\,\rho^2/\tau^2<1.1\cdot10^6\rho^2/\tau^2
\le n$.
\end{proof}

\begin{corollary}[Random formulas]\label{cor:random}
Let $0<\tau<1$, $C=10^9/\tau$ and $n\ge10^9/\tau^2$. In each of the following
cases the formula satisfies, with probability at least $1-1/n$, the conditions
of Corollary~\ref{cor:mean} with $K=9$ and $g=\tau^{3/2}/12$, whatever its unit
clauses are.
\begin{enumerate}
 \item[(a)] Model~(a) with $p\le\frac12$ and $d\ge9C\log n$, for every
 $\varepsilon\in[-1,1]$; and model~(a) with $p=1$ if
 $n(1-\varepsilon^2)\ge5C\log n$.
 \item[(b)] Model~(b) with $q(1+|\varepsilon|)\le\frac12$ and
 $d\ge34C\log n$, for every $\varepsilon\in[-1,1]$.
\end{enumerate}
On such a formula the algorithm of Corollary~\ref{cor:mean} solves Max 2-SAT
in polynomial space and time $\Ostar(c^n)$, with
$c=2^{\max\{\tau,\,1-c_*/\ln2\}}<2$ and $c_*\ge\tau^3/(5.5\cdot10^8)$. The
algorithm is exact on every input; only its running time uses the
probability.
\end{corollary}

\begin{proof}
All pairs have the same distribution, so Proposition~\ref{prop:random}(i)
applies with $\rho=1$ and $\nu=\operatorname{Var}J_{uv}$. In model~(a),
$B=2$ and $\sigma_0=|\varepsilon|p\le p$. If $p\le\frac12$, then $\nu\ge p/2$
and $n\nu\ge d/2\ge C(4\log n+1)$. If $p=1$, then $\nu=1-\varepsilon^2$ and
$\sigma_0\le1$. In model~(b), $B=4$, $\sigma_0=4q|\varepsilon|\le2$ and
$n\nu\ge2qn\ge d/2\ge C(16\log n+4)$. The running time is
Theorem~\ref{thm:main} with $\alpha=81/\sqrt\tau$ and
$c_*=\tau^3/(144\cdot64\cdot81^2\cdot9)$.
\end{proof}

\begin{remark}[Constants and density]\label{rem:random}
No attempt was made to optimize $C$ and $n_0$, and $c_*$ is tiny. The proved
values $K=9$ and $g=\tau^{3/2}/12\approx0.029$ (for $\tau=\frac12$) are not far
from the measured ones in Table~\ref{tab:membership}, which
\texttt{explore\_count\_guessing.py membership} produces. The condition
$d\ge C\log n$ is needed for~\eqref{R1}: $\|R\|$ is at least the largest norm
of a row or a column of $R$, which is of the order of the square root of the
largest number of neighbours of a variable. For $d=o(\log n)$ some variables
have far more than $d$ neighbours, so $K$ is unbounded. Densities between
$\omega(1)$ and $C\log n$ would need a separate treatment of these few
variables, for instance branching on them first.
\end{remark}

\begin{table}[ht]
\centering
\begin{tabular}{lrrrrr}
\toprule
model & $n$ & $\varepsilon$ & $K$ & $g$ & \eqref{R3} ratio\\
\midrule
(a), $p=1$ & 2000 & 0 & 2.42 & 0.034 & 3450\\
(a), $p=1$ & 2000 & 0.5 & 2.44 & 0.034 & 0.66\\
(a), $p=1$ & 2000 & 0.9 & 2.46 & 0.034 & 0.18\\
(a), $d=5$ & 4000 & 1 & 3.45 & 0.028 & 12.5\\
(a), $d=20$ & 4000 & 0.5 & 2.62 & 0.033 & 14.7\\
(a), $d=80$ & 4000 & 1 & 2.47 & 0.034 & 3.74\\
(b), $q=\frac14$ & 2000 & 0 & 2.39 & 0.034 & 1380\\
(b), $q=\frac14$ & 2000 & 0.5 & 2.42 & 0.034 & 0.63\\
(b), $d=20$ & 4000 & 0 & 2.68 & 0.032 & 3360\\
(b), $d=20$ & 4000 & 1 & 2.61 & 0.033 & 7.26\\
\bottomrule
\end{tabular}
\caption{Measured constants of random formulas for the mean-centred
decomposition, $\tau=\frac12$, one formula per row: $K=\|R\|^2/D$, the largest
$g$ allowed by~\eqref{R2}, and the ratio of the right side of~\eqref{R3} to
$\sigma$ (at least $1$ is needed). Proposition~\ref{prop:random} proves $K=9$
and $g\approx0.029$. For a dense, strongly biased formula~\eqref{R3} needs a
larger $n$; the ratio grows like $\sqrt{n(1-\varepsilon^2)}/|\varepsilon|$.}
\label{tab:membership}
\end{table}

\section{Experiments}\label{sec:experiments}

\texttt{explore\_count\_guessing.py exact} runs the algorithm on 60 formulas
with $6\le n\le11$. They are complete parity-biased formulas, planted
bisections, and random formulas with an arbitrary two-group decomposition. In
every case the algorithm found the optimum, with $T=n/2$, and it evaluated
$12{,}375$ of the $46{,}400$ assignments. \texttt{scaling} estimates the
fraction of all assignments that early enumeration evaluates at $T=n/2$, by
sampling (Table~\ref{tab:pass}).

\begin{table}[ht]
\centering
\begin{tabular}{lrrrr}
\toprule
 & $n=60$ & $120$ & $240$ & $480$\\
\midrule
parity bias $\varepsilon=0.9$, plain & $6.3\cdot10^{-4}$ & $1.1\cdot10^{-3}$ &
 $6\cdot10^{-5}$ & $3.3\cdot10^{-4}$\\
\quad with count guessing & $1.3\cdot10^{-4}$ & $0$ & $0$ & $0$\\
parity bias $\varepsilon=1$, plain & $6.3\cdot10^{-4}$ & $1.4\cdot10^{-3}$ &
 $2.5\cdot10^{-4}$ & $8.3\cdot10^{-4}$\\
\quad with count guessing & $0$ & $0$ & $0$ & $0$\\
planted bisection, $\beta=0.9$, plain & $5.6\cdot10^{-3}$ & $9.4\cdot10^{-4}$
 & $2.5\cdot10^{-4}$ & $0$\\
\quad with count guessing & $7.5\cdot10^{-4}$ & $0$ & $0$ & $0$\\
\bottomrule
\end{tabular}
\caption{Fraction of the $2^n$ assignments evaluated by early enumeration at
$T=n/2$; $16{,}000$ samples ($6{,}000$ for $n=480$), so $0$ means below about
$10^{-4}$. One formula per entry.}\label{tab:pass}
\end{table}

With bias $\varepsilon\le0.8$ plain early enumeration also falls below the
sampling resolution from $n=240$ on, consistent with a floor of size
$e^{-\Theta(1/\varepsilon^2)}$. A heuristic for this floor: on a constant
fraction of prefixes, the creators favour one value by $\Theta(\varepsilon n)$.
The count term then adds to $Q$ a product of this imbalance and the imbalance
of the completion, which is as large as the shortfall with constant
probability. The bound of~\cite{CC} for plain early enumeration
degenerates on these formulas, because the aligned part has norm
$|\beta|n\gg\sqrt D$. Count guessing removes the count term from $Q$
altogether.

\section{Relation to the note on group counts}\label{sec:relation}

\begin{itemize}
 \item \emph{Group counts.} The solver of~\cite[Section~5]{RG} exploits exact
 repetition across one cut, with arbitrary structure on both sides, by
 enumerating the sides separately. Count guessing exploits aligned terms that
 are global but noisy, and hands the rest to pruning. The two combine: a
 run of Theorem~\ref{thm:main} may end in the group-count solver instead of
 enumeration whenever its remaining formula has repeated profiles.
 \item \emph{Useful constraints.} In a run, the learned constraints are those
 of $F'_c$. They differ from those of $F$ exactly by the count term, which is
 what made them redundant. Guessing the counts is therefore a way of
 \emph{making} constraints useful, not only of choosing among them.
 \item \emph{One-literal-neighbourhood constraints.} The constraint
 $A_\ell\ge G_\ell+1$ of~\cite[Section~2]{RG} excludes a fraction
 $\Prb[\mathrm{Bin}(d_2(\ell),\frac12)\le G_\ell]$ of the completions of its
 support. This fraction is non-negligible only when $G_\ell\ge
 d_2(\ell)/2-O(\sqrt{d_2(\ell)})$, that is, for strongly biased variables. In
 the regime of Theorem~\ref{thm:main} the combined constraint is the one that
 carries the shortfall.
 \item \emph{A by-product for near twins.} If every variable has an
 $r$-near-twin in the sense of~\cite{NT}, then a spanning star forest of the
 near-twin graph yields at least $n/3$ disjoint $(2r+1)$-near-twin pairs:
 two leaves of a star differ at most at the centre and at their own
 discrepancies. So~\cite[Theorem~3]{NT} gives time
 $\Ostar((2(3/4)^{1/(6r+6)})^n)$ for this natural class, at every density.
\end{itemize}

\paragraph{What is not covered.} Aligned structure of high rank: a blow-up
with groups of size $L$ needs $n/L$ group counts. The resulting
$(L+1)^{n/L}=2^{n\log_2(L+1)/L}$ runs are affordable only when
$\log_2(L+1)/L$ is below the saving of a run, while for bounded $L$ the plain
criterion of~\cite{CC} applies with $K$ of order $L$. Between the two, and for aligned parts that are not
low-rank at all, a different exit is needed. The decomposition must also be
known, except in the mean-centred case. Random formulas are covered only at
density $d\ge C\log n$ (Remark~\ref{rem:random}). Finally, an adaptive version would
measure the runs with the soft count~\cite{SC} rather than
enumerate them.

\begin{thebibliography}{10}
\small\raggedright
\bibitem{BvH16} A.~S.~Bandeira and R.~van Handel. Sharp nonasymptotic bounds on
 the norm of random matrices with independent entries. \emph{Ann. Probab.}
 44(4):2479--2506, 2016.
\bibitem{BLM13} S.~Boucheron, G.~Lugosi and P.~Massart. \emph{Concentration
 Inequalities: A Nonasymptotic Theory of Independence}. Oxford University
 Press, 2013.
\bibitem{CA} \emph{A checkable replacement for the gap-profile condition},
 addendum, 8 October 2026 (\texttt{max2sat\_counting\_addendum.tex}).
\bibitem{CC} \emph{Replacing random signs by neighborhood counts}, working note,
 8 October 2026 (\texttt{max2sat\_counting\_conditions.tex}).
\bibitem{EE} \emph{Max 2-SAT by early constraint-guided enumeration}, working
 note, 8 October 2026 (\texttt{max2sat\_early\_enumeration.tex}).
\bibitem{Mau03} A.~Maurer. A bound on the deviation probability for sums of
 non-negative random variables. \emph{J. Inequal. Pure Appl. Math.} 4(1),
 Article~15, 2003.
\bibitem{NT} \emph{Near-twin neighborhoods as a structural exit}, working note,
 8 October 2026 (\texttt{max2sat\_near\_twin\_pairs.tex}).
\bibitem{RG} \emph{Useful constraints or compression by group counts}, working
 note, 10 October 2026 (\texttt{max2sat\_redundancy\_and\_group\_counts.tex}).
\bibitem{SC} \emph{The soft-count measure for adaptive branching}, working
 note, 8 October 2026 (\texttt{max2sat\_soft\_count.tex}).
\bibitem{Tal95} M.~Talagrand. Concentration of measure and isoperimetric
 inequalities in product spaces. \emph{Publ. Math. IH\'ES} 81:73--205, 1995.
\end{thebibliography}
\end{document}
